\chapter{Replanejamento quíntico com planeza diferencial (Hermite)}
\label{cap:hermite}

\section{Lei de controle}
Arquivos \code{control/hermite\_flatness\_controller.cpp} e \code{hermite\_spline.cpp}.
O quadricóptero é \emph{diferencialmente plano}: posição e yaw (e suas derivadas)
determinam atitude, empuxo e taxas. Por isso basta planejar uma curva suave de
posição e enviar $(\vect v,\vect a,\vect j)$ como \emph{feedforward}.

A cada ciclo constrói-se um polinômio de grau 5 por eixo, $\vect p(t)=\sum_{i=0}^5
\vect c_i t^i$, com 6 condições de contorno:
\begin{align}
  t=0:&\quad \vect p=\vect 0,\ \vect v=\vect v_{ref},\ \vect a=\vect a_{prev}
       \quad\text{(continuidade $C^2$ com o plano anterior)}\\
  t=T:&\quad \vect p=\text{standoff}(\rr+\vect v_tT),\ \vect v=\vect v_t,\ \vect a=\vect 0 .
\end{align}
A duração $T$ é a única variável livre e é otimizada por \textbf{seção áurea}:
\begin{equation}
  \min_{T\in[T_{min},T_{max}]}\ w_T\,T + w_j\!\int_0^T\!\lVert\dddot{\vect p}\rVert^2dt
  + \frac{w_{lim}}{M}\sum_{i=1}^{M}\left[(\lVert\vect v(t_i)\rVert^2-v_{max}^2)_+^2
  + (\lVert\vect a(t_i)\rVert^2-a_{max}^2)_+^2\right].
\end{equation}
O comando é o plano avaliado em $t=\Ts$.

\section{Interpretação dos pesos}
\begin{itemize}
  \item $w_T$ (\code{weight\_time}): ``pressa''. Maior $w_T\Rightarrow T$ menor e manobra
        mais agressiva.
  \item $w_j$ (\code{weight\_jerk}): suavidade. Para um deslocamento $\Delta$ em
        repouso-repouso, o jerk mínimo vale $\int j^2 = 720\,\Delta^2/T^5$. O equilíbrio
        $w_T = 5\cdot720\,w_j\Delta^2/T^6$ dá
        $T^\star\approx(3600\,w_j\Delta^2/w_T)^{1/6}$. Com $w_j/w_T=0{,}02$ e
        $\Delta=\SI{3}{m}$: $T^\star\approx\SI{2,9}{s}$.
  \item $w_{lim}$: penalidade suave dos limites (quanto maior, mais próximo de restrição
        dura).
\end{itemize}

\section{Equivalência linear}
Para pequenos erros a lei também é aproximadamente linear no estado. Avaliando a
derivada do plano em $t=\Ts\ll T$, o termo dominante dá uma dinâmica de erro com
constante de tempo $\propto T^\star$. Calcule o ganho equivalente numericamente
(perturbando $e_p$ e $e_v$ e lendo $\Delta\vect a$) em \code{projeto\_hermite.m} e
use as mesmas margens dos outros capítulos. Atenção: como $T^\star$ depende de
$\Delta$, a ``banda'' muda com o tamanho do erro (o controlador é não linear).

\section{Ajuste ótimo}
Os parâmetros $(w_T, w_j)$ entram em \code{otimizar\_ganhos.m} como qualquer outro
$\theta$. Fixe $w_{lim}$ alto (50--100) e $T_{min}\ge 4\Ts$.
